\section{Discussion}\label{sec:discussion}

Statements in this section are labelled \emph{proved} (a merged record with nonauthor review), \emph{read} (read in bounded nonauthor slices from the same account), \emph{author-side}, \emph{numerical} (computation or simulation without certification) or \emph{expectation} (not proved); conditional scopes are stated in place. None is used in \S\ref{sec:proof}.

\subsection{Why \texorpdfstring{$-1/3$}{-1/3}, and what depends on the law}\label{ssec:why}
In the separation variable $r=|M-S|$ and the height gap $h$, the ordered maximum/saddle pair measure is the polar volume $r^{d-1}$ times the two-site pin density, of order $r^{-(d+3)}$, times the determinant scale $r^2$; that is, $r^{-2}\,dr\,dh$ in every dimension. A pinned pair has gap $h=kr^3$, so $dh=r^3\,dk$ at fixed $r$ and the measure is $r\,dr\,dk$; this is \eqref{eq:radial} with $k=\kappa/6$. Pushing forward to the lifetime $\ell=kr^3$ at fixed $k$, $r\,dr=\tfrac13k^{-2/3}\ell^{-1/3}\,d\ell$ (cf.\ \eqref{eq:rdr}), which gives $\ell^{-1/3}$ with no dependence on $d$; the corresponding theorem of \rec{parent} (Theorem~C) holds for every $d\ge2$, and the present $d=3$, side-24 statement is one instance. The exponent is the Whitney-fold exponent: in a one-parameter family a maximum/saddle pair annihilates through $f_t(x)=x^3/3-(t_0-t)x$, with separation proportional to $(t_0-t)^{1/2}$ and gap proportional to $(t_0-t)^{3/2}$, so the time spent with gap in $d\ell$ is proportional to $\ell^{-1/3}\,d\ell$. The Gaussian law enters only the constant: by \rec{parent}, (15.2), $c_{d,L}$ is an angular average of $p_G(0)\,p_{V_u}(0)\,\tau_u^{4/3}D_u$, that is, of densities and of a cone moment of the $3$-jet at the collision point, and it scales as $a^{-2/3}$ under $f\mapsto af$. The proof separates accordingly: the elder rule near a collision (the cap theorem), the barrier $\lambda_{\min}(-D^2f(M))\le(\ell K^2/c_L)^{1/3}$ for a bar of lifetime $\ell$, and the pinned-pair identities are deterministic statements about $C^3$ functions; the Gaussian structure supplies the regression couplings and the uniform moment and density bounds. We therefore expect, but do not prove, that the exponent is shared by stationary fields with $C^7$ paths, finite moments of the derivative suprema and a continuous positive density of the $3$-jet at a point, with $c$ given by the analogue of (15.2) of \rec{parent} (\emph{expectation}).

The near-diagonal law has a second scale. The leading term comes from separations $r\asymp\ell^{1/3}$; at $r\asymp\ell^{1/4}$ the cubic coefficient is comparable to the separation ($\ell\asymp r^4$, the cusp organizing the fold family), and there the first correction arises: on that scale the pinned landscape is a universal quartic ridge on which a maximum/saddle pair of the right types is an elder pair in only one third of its type window, so the local contact law overcounts and the correction $c_1\ell^{1/4}$ of Theorem~\ref{thm:three-term} is negative.

\subsection{The corrections}\label{ssec:corrections}
At separations $r\asymp\ell^{1/4}$ the rescaled pinned landscape converges to a universal random quartic whose superlevel topology reduces exactly to the one-dimensional ridge
\[
g=\vartheta\bigl[2(X+\tfrac12)^2(X-1)+3\varphi(X^2-\tfrac14)^2\bigr];
\]
where $\vartheta=\ell/r^4$ is the cusp variable (the records \rec{chain} and \rec{soft} write it $\kappa$; here $\kappa$ keeps its meaning from \eqref{eq:pins}, so $\vartheta=\kappa/(6r)$); the pair is typed if and only if $|\varphi|<1$, and the saddle is the maximum's elder partner if and only if $|\varphi|<1/3$ (at $\varphi=\pm1/3$ two competing saddles, respectively maxima, have equal heights). Integrating the resulting loss over the scale gives \eqref{eq:c1} (\emph{proved}, \rec{chain}). Through the order $\ell^{1/3}$, the elder rule enters the expansion only through this cusp scale: there the decision is stable with an explicit margin proportional to $\bigl||\varphi|-1/3\bigr|$, while at the fold scale and at larger separations the elder rule costs a power of $\ell$ beyond $\ell^{1/3}$. So the third coefficient of the elder law is the same fold-scale constant $c_2$ as for the candidate law (\emph{proved}, \rec{chain}). Fold-scale rejection does enter at the next order, through the term $R_{2/3}\ell^{2/3}$ of Remark~\ref{rem:beyond}.

At large $\vartheta$ the rejected pairs are soft: exactly one transverse curvature is small, of order $\gamma_1^2/\vartheta$, and the rejected cusp kernel decays like $F_0/\vartheta$ in every dimension; for the Gaussian kernel in $d=3$, $\iint F_0=125\sqrt{30}/(192\pi^3)\approx0.1150$ (\emph{proved}, \rec{soft}). The fold-scale rejection rate is the selection-loss coefficient of Remark~\ref{rem:mechanism}; for the Gaussian kernel it equals $F_0H(k)$ with $H(k)=1+\tfrac{12}{25}k^2+h_{7/2}k^{7/2}+O(k^4)$ and an explicit $h_{7/2}<0$ (\emph{proved} at the conditional scope of \rec{soft}, Theorem~B). These soft rejected pairs produce the $\ell^{2/3}$ term $R_{2/3}$ of the rejected density in Remark~\ref{rem:beyond}; in a full-field simulation in $d=3$, rejected adjacent pairs with $r/\ell^{1/4}<0.8$ number 11, against $10.7$ predicted and $1{,}329$ from the cusp kernel alone (\emph{numerical}, \recs{mc}{soft}).

\emph{Size of the corrections.} By \eqref{eq:relative}, the total relative correction of the three-term law is about $-1.5\%$ at $\ell=10^{-4}$, $-5\%$ at $10^{-3}$ and $-17\%$ at $10^{-2}$; the two corrections partially cancel. A direct simulation of the elder-rule persistence of the periodized Gaussian field (sides $64$ in $d=2$ and $16$ in $d=3$), with exact critical values and exact one-dimensional connectivity by ascending lines, agrees with the three-term law over $10^{-4}\le\ell\le10^{-2}$ (data over law $1.006\pm0.004$ in $d=2$ and $1.008\pm0.010$ in $d=3$) and rejects the leading and two-term laws; it also shows the cusp-scale elder window $|\varphi|<1/3$ directly. Adding the $\ell^{2/3}$ term of Remark~\ref{rem:beyond} improves the fit further ($\chi^2$ from $84$ to $23$ in $d=3$ and from $44$ to $7$ in $d=2$, over $14$ bins on $[10^{-4},3\cdot10^{-2}]$), leaving a small negative next-order term (\emph{numerical}, \rec{mc}). The simulation illustrates the law; it does not prove it, and its residuals are not evidence for an $\ell^{1/2}$ term.

\subsection{Selection loss in dimension three}\label{ssec:d3selection}
In $d\ge3$ the hard transverse directions do not change the local decision: maximizing over the hard fibre reduces the decision window to the planar model up to an error of order $r$ in $C^2$, so the planar elder/rejected certificate transfers to the actual field with margins of the same order (\emph{author-side}: addenda A3 and A3.1 with their erratum, read in bounded nonauthor slices, \rec{qs3}). On compact marks the exhaustion has been carried out in $d=3$: the failure-measure rate holds at every $\beta<1$, with
\[
1-p_r=r^3\alpha^{(3)}+O\bigl(r^4\log(1/r)^7\bigr),
\]
where $\alpha^{(3)}$ is the local cluster mass of Remark~\ref{rem:mechanism}, identified with the fold-scale coefficient $F/(kA_0)$ of \rec{mech}. As a corollary, on compact windows of births and gaps the rejected density is $C^{(3)}_{\mathrm{fail}}\ell^{2/3}+O(\ell\log(1/\ell)^7)$, with $C^{(3)}_{\mathrm{fail}}=\int A_0\,\alpha^{(3)}/(3k^{5/3})\,db\,dk\,d\sigma(u)$ over the window and $A_0$ the pinned-pair weight of \rec{mech} (\emph{read}, and conditional on the author-side chain A3.3--A3.8; independence credit~0; \rec{qs3}). All constants are existential, and no rate of this type is claimed for the unrestricted population.

\subsection{Dimension one}\label{ssec:d1}
In dimension one every ingredient of the short-bar expansion can be made explicit. For a smooth stationary Gaussian process on a circle, and on the line under a mild spectral hypothesis, the persistence lifetime density satisfies
\[
\nu(\ell)=C_0\ell^{-1/3}+C_1\ell^{1/4}+2B_2\ell^{1/3}+O(\ell^{3/4}),
\]
with $C_1<0$ and $B_2$ explicit in the spectral moments; $B_2>0$ when $\lambda_2\lambda_6\le25\lambda_4^2$, for instance for the Gaussian kernel, but its sign is not fixed in general (\emph{proved} on the circle, Theorems 1D and 1D$^+$; on the line \emph{author-side}, read by four nonauthor slices and merged; \rec{d1}). The coefficient $C_1$ is the one-dimensional value of the cusp formula \eqref{eq:c1}, the elder window of the cusp ridge is the same $|\varphi|<1/3$, and the remainder is of the order of the first correction to the cusp term. For the covariance $e^{-x^2/2}$ the crest-to-trough amplitude density per crest is $0.19971814\,h^{-1/3}-0.27165891\,h^{1/4}+0.41725471\,h^{1/3}+O(h^{3/4})$. A simulation over $8.4\cdot10^7$ correlation lengths reproduces the three-term law to $0.3\%$ and the coefficient of the rejected pairs to $1\%$ (\emph{numerical}). For the Gaussian kernel the second- and third-order terms of the lifetime density nearly cancel for $10^{-4}\le\ell\le10^{-2}$ (relative coefficients $C_1/C_0=-2.07$ and $2B_2/C_0=+2.09$), a caution against reading corrections off histograms at moderate $\ell$.

\subsection{The plane}\label{ssec:plane}
In dimension two the local pairing rule is established for the actual Gaussian field with a rate. On the bounded soft layer, the $H_0$ pairing of the designated pin pair agrees with the model's elder decision except on a set of weighted probability $O(r^{7/2})$, against a layer mass of order $r^3$ (\emph{proved}, \rec{planar}). On compact gap intervals the selection loss satisfies $1-p_r=r^3(\alpha_1+\alpha_2)+O\bigl(r^4\log(1/r)^4\bigr)$, so that on compact windows of births and gaps the rejected density is $C_{\mathrm{fail}}\ell^{2/3}+O(\ell\log(1/\ell)^4)$ (\emph{read}, and conditional on C103, A4.1 and A4.2; independence credit~0; \rec{planar}). These are statements about $d=2$; Theorem~\ref{thm:main} is a $d=3$ statement.
